% Agentic AI — Final Project Report
% Authored directly in LaTeX. Compile with: xelatex report.tex (twice for TOC).
\documentclass[11pt,a4paper]{article}

% ---------- Packages ----------
\usepackage[margin=1in]{geometry}
\usepackage{fontspec}
\usepackage{microtype}
\usepackage{setspace}
\setstretch{1.15}

\usepackage{xcolor}
\definecolor{accent}{HTML}{2C5282}
\definecolor{muted}{HTML}{4A5568}
\definecolor{codebg}{HTML}{F7FAFC}
\definecolor{codeborder}{HTML}{E2E8F0}
\definecolor{tablehd}{HTML}{EDF2F7}

\usepackage{booktabs}
\usepackage{array}
\usepackage{tabularx}
\usepackage{longtable}
\renewcommand{\arraystretch}{1.25}

% enumitem unavailable on this TeX install; default itemize spacing is fine

\usepackage{listings}
\lstset{
  basicstyle=\ttfamily\footnotesize,
  backgroundcolor=\color{codebg},
  frame=single,
  rulecolor=\color{codeborder},
  framerule=0.4pt,
  framexleftmargin=4pt,
  framexrightmargin=4pt,
  framextopmargin=4pt,
  framexbottommargin=4pt,
  xleftmargin=0pt,
  breaklines=true,
  showstringspaces=false,
  columns=flexible,
  aboveskip=0.8em,
  belowskip=0.8em,
}

\usepackage[hidelinks,colorlinks=true,linkcolor=accent,urlcolor=accent]{hyperref}

% Both titlesec and sectsty unavailable on the basic TeX install. Override
% \section / \subsection directly so headings still get the accent colour.
\makeatletter
\renewcommand\section{\@startsection{section}{1}{\z@}%
   {-1.6ex \@plus -0.5ex \@minus -0.2ex}%
   {0.8ex \@plus 0.2ex}%
   {\Large\bfseries\color{accent}}}
\renewcommand\subsection{\@startsection{subsection}{2}{\z@}%
   {-1.1ex \@plus -0.4ex \@minus -0.2ex}%
   {0.6ex \@plus 0.2ex}%
   {\large\bfseries\color{accent!85!black}}}
\makeatother

\usepackage{fancyhdr}
\pagestyle{fancy}
\fancyhf{}
\fancyfoot[C]{\small\color{muted}\thepage}
\renewcommand{\headrulewidth}{0pt}

% Use a clean sans-serif throughout. Charter has good readability in body text;
% fall back to system fonts that have unicode arrows.
\setmainfont{Charter}[Ligatures=TeX]
\setmonofont{Menlo}[Scale=0.9]

% Convenience macros
\newcommand{\rarr}{\(\rightarrow\)}
\newcommand{\code}[1]{\texttt{\small #1}}

% ---------- Document ----------
\begin{document}

% ===================== Title block =====================
\begin{center}
\vspace*{-1.2em}
{\Huge\bfseries\color{accent} Agentic AI --- Final Project Report\par}
\vspace{1.2em}
\renewcommand{\arraystretch}{1.1}
\begin{tabular}{l@{\hspace{2em}}l}
\textbf{Muhammad Taha}    & 22i-0983 \\
\textbf{Ibtehaj Haider}   & 22i-0767 \\
\textbf{Hussain Waseem}   & 22i-0893 \\
\end{tabular}
\renewcommand{\arraystretch}{1.25}
\end{center}

\vspace{1.5em}
\hrule
\vspace{1.5em}

% ===================== 1. Executive summary =====================
\section{Executive summary}

This project implements the complete \emph{prompt \rarr{} polished short film}
pipeline specified in the course brief. A single English-language prompt produces,
without further human creative intervention, a finished animated short: structured
script, voiced dialogue, per-scene mood-music, background visuals, subtitles,
per-character portraits, and a final composited MP4. An intelligent edit agent
then accepts free-text commands ("make scene 2 darker", "speed up this scene",
"regenerate the script"), routes each to the minimal affected phase, and writes
a versioned snapshot so any edit is reversible with branching semantics.

\smallskip
\textbf{Delivered in this submission:}
\begin{itemize}
  \item All five phases implemented and integration-tested.
  \item 47 passing unit + integration tests covering the shared schema, MCP tool
        layer, state manager (with branching), intent classifier (17 query types),
        planner, subtitle timing, and backend routes.
  \item An end-to-end demo run (3-scene Neo-Tokyo detective vignette) renders in
        \textbf{$\sim$5 minutes} cold-start on an Apple M4 Air, with 3 edit actions
        + 2 undos applied in under 30 seconds each.
  \item A React web UI serving live phase progress via WebSocket, scene thumbnails,
        version history with active-pointer + lineage tooltips, and the edit-chat
        interface.
  \item \textbf{Wav2Lip-driven lip-sync} on speaker close-ups via a
        mouth-region alpha-blend that keeps the rest of the frame at native
        portrait resolution.
  \item Zero paid APIs are required: Groq (free tier), Pollinations (keyless),
        edge-tts (keyless), \code{ffmpeg} (local), Wav2Lip (local CPU). Optional
        drop-in upgrades are wired for ElevenLabs (voice quality) and fal.\,ai
        (HD image gen) when keys are configured in \code{.env}.
\end{itemize}

% ===================== 2. System architecture =====================
\section{System architecture}

\subsection{Five modules, one JSON contract}
The system is five loosely-coupled modules connected by a single Pydantic state
object, \code{PipelineState}. Nothing else crosses module boundaries --- the
edit agent, web UI, and any future module that wants to participate just reads
and writes that one object.

\smallskip
\begin{lstlisting}[language={}]
   prompt
     |
     v
  Phase 1 (story_agent)        --->  StoryState
     |
     v
  Phase 2 (audio_agent)        --->  TimingManifest
     |
     v
  Phase 3 (video_agent)        --->  VideoOutput  --->  final_*.mp4
     ^
     |  re-enter any phase
     |
  Phase 5 (edit_agent)  <-->  StateManager (SQLite log + asset freeze)
     ^
     |
  Phase 4 (web UI / FastAPI)
\end{lstlisting}

Each phase exposes both a \emph{full-run} entry point (e.g. \code{run\_story\_phase})
and partial entry points used by the edit agent (e.g.
\code{regenerate\_scene\_audio}, \code{regenerate\_scene\_background}). This dual
entry surface is what makes precise per-intent re-runs cheap: \emph{"make scene 2
darker"} regenerates one background and re-stitches the final video; the other
scenes' assets are untouched.

\subsection{Shared JSON schema}
The Pydantic schema in \code{shared/schemas/} is the single source of truth.
Every value crossing phase boundaries is typed:

\begin{itemize}
  \item \code{PipelineState} --- top-level container: \code{job\_id}, \code{prompt},
        \code{num\_scenes}, \code{style}, \code{story}, \code{audio}, \code{video},
        \code{phase\_status}, \code{version}.
  \item \code{StoryState} --- title, logline, language, list of \code{Scene}
        (location, action, characters, dialogue lines with emotion + visual cue),
        list of \code{Character} (name, role, appearance, voice\_style, gender,
        reference\_style, image\_path).
  \item \code{TimingManifest} --- per-line \code{AudioSegment}s with start/end
        offsets, plus per-scene merged audio paths and BGM track paths.
  \item \code{VideoOutput} --- per-scene \code{SceneClip}s (background path, raw
        ken-burns clip, composed scene MP4) plus the final stitched MP4.
\end{itemize}

\subsection{MCP tool layer}
A small framework lives under \code{mcp/}: \code{BaseTool} (with a \code{ToolSpec}
declaring name + category + JSON schema), a singleton \code{ToolRegistry}, and
\code{ToolExecutor} which wraps every call with retry, structured event emission,
and automatic \code{job\_id} injection so every tool's outputs route into a
per-job subdirectory. The executor's introspection step
(\code{inspect.signature(tool.run)}) only injects \code{job\_id} when the tool
declares the parameter, so legacy LLM tools that don't care about it remain
untouched.

\bigskip
% ===================== 3. Phase-by-phase =====================
\section{Phase-by-phase implementation}

\subsection{Phase 1 --- Story, script \& character design (15\%)}

\textbf{Inputs:} free-form prompt + scene count + style hint.
\textbf{Outputs:} \code{StoryState} with scenes, dialogue, character roster,
character portrait images.

The story agent makes two LLM calls via Groq's free tier (Llama 3.3 70B in JSON
mode): \code{generate\_story} produces title, logline, scenes with dialogue
turns, and \code{design\_characters} normalises a roster with appearance,
\code{voice\_style}, \code{gender}, and \code{reference\_style}. Strict JSON-mode
output plus a Pydantic post-validation step removes the entire class of
"the LLM almost gave us valid JSON" bugs.

\textbf{Prompt-engineering iterations.} An early failure mode was the LLM
populating the \code{line} field with \emph{stage directions} ("Rain pours
down", "Shadows hide faces") instead of spoken dialogue. We tightened the
system prompt with explicit GOOD vs. BAD examples and a rule that lines must be
8--25 words, contain contractions / hesitation, and never describe what the
camera sees (that goes in \code{visual\_cue}).

A second, subtler variant survived the system-prompt fix: parenthetical and
markdown stage directions \emph{embedded inside} otherwise-spoken lines
(\texttt{"(whispers) I see something"}, \texttt{"*sighs* fine"},
\texttt{"[shouting] now!"}). The TTS provider read these aloud verbatim
("whispers I see something") because nothing stripped them on the way to
synthesis. Resolved with a regex pass in \code{normalise\_story} that strips
\texttt{(\ldots)}, \texttt{[\ldots]}, and \texttt{*\ldots*} segments before the
line lands in the \code{DialogueTurn}. The cleaned line is used both for TTS
and for on-screen subtitles, so audio and captions stay in lockstep.

\textbf{Style-aware prompting.} The user-facing style dropdown
(\code{cinematic / photorealistic / anime / noir / cyberpunk / watercolor})
originally fed a single token into the image prompt while the surrounding
prompt hardcoded \emph{"photorealistic, ultra detailed, sharp focus, 8k"} as
quality tail. The contradictory tokens caused diffusion models to ignore the
selected style and resolve toward photorealism every time. Replaced the
hardcoded tail with a per-style token map: \code{anime} now anchors with
\texttt{"anime illustration: \ldots cel-shaded, vibrant flat colors, manga
aesthetic, no photorealism"}, \code{noir} with \texttt{"film noir frame: \ldots
black and white, hard shadows, 1940s detective film"}, \code{cyberpunk} with
\texttt{"cyberpunk concept art: \ldots neon-soaked, blade runner aesthetic"},
etc. Cinematic and photorealistic remain distinct ("35\,mm film grain,
anamorphic, color graded" vs.\ "hyperrealistic photograph, dslr, natural
lighting"). The cache key includes the style token, so existing jobs cache-miss
into the corrected prompt on next render.

\textbf{Character portraits} are fetched from Pollinations (square 1024 px
images), with a Pillow-rendered placeholder fallback so the pipeline never
hard-fails on the "output completeness" rubric criterion.

\subsection{Phase 2 --- Audio generation \& integration (15\%)}

\textbf{Inputs:} \code{StoryState}.
\textbf{Outputs:} per-line WAV files, per-scene merged audio, \code{TimingManifest}
with millisecond offsets.

\textbf{Voice assignment.} Each character is hashed (MD5 of name) into a
gender-aware voice pool. Hash determinism guarantees the same character keeps
the same voice across all scenes. The pool is split by gender; an explicit
character.gender field (populated by the LLM character designer) selects which
half to hash into. We use Microsoft \emph{Multilingual Neural} edge-tts voices
by default --- the highest-quality keyless tier --- and ElevenLabs Multilingual
v2 when \code{ELEVENLABS\_API\_KEY} is set in \code{.env}, with a graceful
fallback if the EL quota is exhausted.

\textbf{Emotion playbook.} Each scene's dialogue line carries an \code{emotion}
attribute (extracted by the LLM). At synthesis time we apply both a
\emph{text-cue} rewrite (\texttt{"!!"} for shout, \texttt{"..."} for sad,
lowercase + ellipsis for reflective) and an emotion-specific FFmpeg post-filter:

\begin{tabularx}{\textwidth}{l X X}
\toprule
\textbf{Emotion} & \textbf{Text rewrite} & \textbf{FFmpeg post-filter} \\
\midrule
whisper & --- (text unchanged) & \code{volume=0.45, highpass=f=100, treble=g=4} \\
shout & strip terminal \texttt{.}, append \texttt{!!} & \code{volume=1.35, acompressor, treble=g=2} \\
angry & append \texttt{!} & \code{acompressor, treble=g=2} \\
sad & append \texttt{...} & \code{atempo=0.93, lowpass=f=4500, volume=0.88} \\
happy / excited & append \texttt{!} & \code{atempo=1.05, treble=g=2} \\
panic / fear & append \texttt{!} & \code{atempo=1.08, treble=g=3, acompressor} \\
\bottomrule
\end{tabularx}

The on-screen subtitle still uses the original line text, so captions stay
clean regardless of the cued text.

\textbf{BGM} is mood-driven: each scene's mood maps to a track in the
\emph{Ambient Film Music} CC-BY 3.0 library (Serge Quadrado, sourced from
archive.org via \code{scripts/fetch\_bgm.py}). If a track is missing, a
multi-oscillator chord pad synthesised in pure FFmpeg \code{lavfi} is used as
fallback. The dialogue \rarr{} BGM mix uses a static $-7$\,dB bed instead of
sidechain ducking, which (despite its theoretical appeal) suppressed the BGM
into inaudibility against TTS that already varies in level.

\subsection{Phase 3 --- Video generation \& composition (20\%)}

\textbf{Inputs:} \code{StoryState} + \code{TimingManifest}.
\textbf{Outputs:} per-scene MP4s and the final stitched \code{final\_<job\_id>.mp4}.

\textbf{Per-line close-up flow.} Rather than a single Ken-Burns over the scene
background, each scene is rendered as: a 1.8\,s \emph{establishing wide shot}
of the location (Ken-Burns over the BG, brief mood-BGM under it), then a
sequence of \emph{speaker close-ups}: each dialogue line uses that speaker's
portrait, top-aligned 16:9 crop (so the face stays in frame), Ken-Burns'd over
the line's audio duration, with the line's subtitle burned on. Hard cuts
between speakers give classic shot/reverse-shot feel.

\textbf{Wav2Lip mouth-region blend.} The close-ups originally showed a static
portrait whose lips never moved while the speaker's voice played, which broke
the illusion of dialogue. We integrated the pretrained
\href{https://github.com/justinjohn0306/Wav2Lip}{Wav2Lip GAN} (vendored under
\code{third\_party/Wav2Lip}, checkpoint \code{wav2lip\_gan.pth}) to synthesise
lip motion from the dialogue WAV. Three engineering hurdles surfaced:

\begin{itemize}
  \item \textbf{NumPy 2.x incompatibility.} The fork uses retired aliases
        \code{np.int} which 2.x removed; patched call sites in
        \code{face\_detection/utils.py} to \code{np.int32}.
  \item \textbf{Bundled face-detector weights missing.} The inference script
        hardcoded \code{checkpoints/mobilenet.pth} (not shipped); patched to
        \code{model\_path=None} so \code{batch\_face} auto-downloads the
        canonical mobilenet RetinaFace weights to the user's torch hub cache
        on first run.
  \item \textbf{Whole-frame softness from the 96$\times$96 face patch.}
        Wav2Lip internally renders the face at 96$\times$96 and upscales it
        to fit the detected face bbox. Used as the full visual it noticeably
        blurs the entire face. We resolved this by treating Wav2Lip's output
        as a \emph{motion source} for the mouth area only: the static portrait
        is rendered at native resolution as the base, an elliptical alpha
        mask is built around the detected mouth landmarks (RetinaFace returns
        mouth-corner points; the ellipse covers them plus a 7\,\% face-height
        chin allowance), and Wav2Lip's clip is alpha-merged onto the base via
        FFmpeg only inside that ellipse. A 4\,\% Gaussian feather softens the
        seam. Outside the ellipse, every pixel is the original portrait at
        native resolution; only the mouth (where lip-motion lives) carries
        Wav2Lip pixels.
\end{itemize}

\textbf{Quality-recovery chain on the masked region.} Even confined to the
mouth ellipse, the 96\,px Wav2Lip patch is perceptibly softer than the
surrounding portrait. The lip-sync stream goes through a 2$\times$ Lanczos
upscale, two-radius unsharp (3\,px and 7\,px), AMD's CAS (Contrast-Adaptive
Sharpening, halo-free by design and intended for upscaled content), and a
Lanczos downscale back to target. The non-mouth pixels never see this filter.
A \code{try/except} around the \code{mouth\_blend\_clip} tool falls back to
plain Ken-Burns motion if face detection fails on a stylised illustration, so
one bad portrait never kills the render.

\textbf{Subtitle burning} uses Pillow-rendered transparent PNGs overlaid via
FFmpeg's \code{overlay=enable='between(t,a,b)'} filter --- the standard
\code{subtitles}/\code{ass} filters require libass which Homebrew's stock
\code{ffmpeg} ships without. The PNG path is libass-free, supports millisecond
timing, and renders with anti-aliased text + drop-shadow on a translucent
rounded background plate so captions remain legible on bright scenes.

\textbf{Final compositor} stitches scene MP4s with FFmpeg \code{xfade}
transitions and \code{acrossfade} on the audio, then applies
\code{-movflags +faststart} so the \code{moov} atom lives at the front of the
file and browsers can begin progressive playback before the full MP4 is
buffered.

\textbf{Encoder choice.} We initially used \code{h264\_videotoolbox} on macOS
for $\sim$6\,$\times$ speed-up over software encoding, but discovered late in
development that videotoolbox silently outputs \code{yuvj420p} (full-range,
JPEG-style colour space) which Firefox 150 refuses to play in HTML5
\code{<video>}. Switched the default encoder to \code{libx264} which produces
clean \code{yuv420p} + \code{color\_range=tv} that every browser handles. The
videotoolbox path remains available via the \code{VIDEO\_ENCODER} env var.

\subsection{Phase 4 --- Web interface (10\%)}

A FastAPI backend with a WebSocket progress stream, paired with a React 18 SPA
served as a pre-built JSX bundle. The bundle is built once with
\code{esbuild}; users don't need to run \code{npm install} for the demo --- the
backend's \code{/static} mount serves the JS directly.

\textbf{Real-time progress.} The pipeline emits structured events
(\code{phase\_start}, \code{tool\_start}, \code{tool\_ok}, \code{scene\_composed},
\code{final\_ready}, \code{snapshot}, \code{job\_complete}) that the WebSocket
fans out to all connected clients of the relevant \code{job\_id}. The UI
renders these as a live event log.

\textbf{Phase rerun} buttons in the sidebar invalidate the affected
\code{PipelineState} subtrees and re-trigger only the requested phase ---
useful when the LLM produces a flat scene and you want to roll the dice again
without losing the audio you already rendered. Manual reruns
\textbf{cascade by default}: clicking \emph{Rerun story} regenerates the
script and then auto-runs audio + video against the new dialogue (since the
existing TTS wavs and scene composites are stale once dialogue lines change).
\emph{Rerun audio} cascades into video. The cascade produces a single snapshot
row in the version history, not three, so the dropdown stays clean. Edit-agent
automatic reruns can still pass \code{force=false} to skip the cascade --- they
know exactly which assets they touched.

\textbf{Force-regenerate cache invalidation.} Every tool in the pipeline
caches by content-hash filename to keep edit roundtrips cheap. Without
intervention, this means a manual rerun returns the previously cached file
byte-for-byte and produces no visible change --- which is correct
deterministic behaviour but bad UX (users expect "rerun" to do work). Manual
reruns now wipe the per-job cache directories before invoking the phase
handler:

\begin{center}
\begin{tabularx}{\linewidth}{l X}
\toprule
\textbf{Phase} & \textbf{Wiped (recursive)} \\
\midrule
\code{video}    & \code{video/<job\_id>/}, \code{scenes/<job\_id>/} (covers ken-burns, wav2lip, mouth-blend, scene backgrounds, alpha masks, scene composites) \\
\code{audio}    & + \code{audio/<job\_id>/} (TTS wavs, BGM tracks) \\
\code{story}    & + \code{images/<job\_id>/} (character portraits) \\
\bottomrule
\end{tabularx}
\end{center}

\noindent The recursion matters: scene backgrounds live two levels deep
(\code{video/<job\_id>/scene\_backgrounds/bg\_*.png}) and a flat wipe leaves
those PNGs intact, masking the regen. With recursive wipe, image gen
re-fetches from Pollinations with a fresh \code{seed=uuid4()} and every
establishing shot visibly differs across runs. The wipe helper lives in
\code{shared/utils/phase\_cache.py} so the rerun route and the edit-agent
executor share one source of truth on what "invalidate phase X" means.

A separate concern is \emph{LLM determinism}: Groq calls were initially
pinned at \code{temperature=0.0} for asset-cache stability during development.
At that setting the same prompt returned byte-identical text every time, so
rerunning the story phase did nothing visible. Bumped to
\code{temperature=0.8} for stories and 0.5 for character rosters, plus a
randomly-chosen \emph{variation directive} appended to the prompt on
force-regen reruns ("Try a fresh angle on this story", etc.) so the LLM
explores different dialogue rhythms even within the same world.

\textbf{Resume picker.} The frontend persists \code{job\_id} to
\code{localStorage} and queries \code{/api/jobs} for a list of past runs, so
closing the browser doesn't lose your active session and you can switch
between jobs from a dropdown.

\textbf{Version-history panel.} Each row shows the version number, an arrow
indicating its parent (\code{$\leftarrow$ vN}), the triggered-by tag, and a
human change summary. The currently \emph{active} version is highlighted with
an accent-coloured border + "$\bullet$ active" badge. Hovering any row reveals
its full lineage chain (e.g.\ \code{lineage: v6 \rarr{} v3 \rarr{} v1}).

\textbf{Cache-bust on snapshot.} The \code{<video>} element appends
\code{?b=<active\_version>} to its \code{src}; whenever \code{refreshAll}
fetches \code{/api/versions/<job\_id>} and discovers the active version has
changed, the bust value updates and the player calls \code{.load()} to drop
its buffer and refetch. Without this, every rerun produced a new MP4 at the
same path on disk but the browser kept showing the old buffered video. The
sync-to-active-version (rather than incrementing on UI actions only) means
\emph{any} new snapshot --- UI-driven, CLI-driven, or programmatic ---
correctly invalidates the player.

\subsection{Phase 5 --- Edit agent + undo (20\%)}

The edit agent is a four-node LangGraph: \code{classify\_intent} \rarr{}
\code{plan\_execution} \rarr{} \code{execute} \rarr{} \code{snapshot}. The
classifier tries Groq first (LLM-mode) for nuanced understanding; on any error
it falls back to a deterministic keyword rule set so the demo never blocks
offline. The planner maps an \code{EditIntent} to a list of (handler, kwargs)
steps; the executor dispatches each step to a story / audio / video module
function.

\textbf{Generic visual modifications.} A single \code{modify\_scene\_visual}
intent handles the full free-text design surface --- "make scene 1 darker",
"more orange", "add buildings", "make it nighttime" --- by extracting four
structured params (\code{prompt\_suffix}, \code{brightness\_delta},
\code{saturation\_delta}, \code{hue\_shift}) from the directive and applying
them: the suffix is appended to the scene-background prompt for re-fetch, and
the deltas drive an FFmpeg \code{eq + hue} colour-grade pass on the rendered
scene MP4.

\textbf{State versioning.} Every snapshot is an append-only row in a SQLite log
(\code{state.sqlite}) with a \code{parent\_version} pointer. A separate
\code{active\_version} table tracks which row is currently checked out.
\textbf{Undo and restore are pointer moves}, not new rows --- they update
\code{active\_version} to a target version and copy that version's frozen
assets back into the working tree. Subsequent edits create new versions whose
parent is whatever the active pointer was at edit time, giving us git-style
branching: undo to v2, edit again, and v4 has \code{parent\_version=2} (not
v3, the abandoned branch).

\bigskip
% ===================== 4. Challenges =====================
\section{Challenges and resolutions}

\begin{description}
  \item[Cross-job filename collision.] Per-scene MP4 paths
    (\code{scenes/scene\_01.mp4}) were shared across jobs, so generating job B
    overwrote job A's working files.  \textbf{Fix:} every output tool now
    routes through a \code{job\_dir(BASE, job\_id)} helper that returns
    \code{<BASE>/<job\_id>/} --- structurally impossible for two jobs to
    overwrite each other.
  \item[ElevenLabs read \texttt{[whispers]} as text.] The Multilingual v2 model
    advertises support for inline audio tags but in practice often pronounces
    them aloud. \textbf{Fix:} dropped the tag prefix entirely and implemented
    whisper acoustically via FFmpeg post-processing (volume cut + high-pass +
    treble lift).
  \item[Edge-tts retired voices.] Several voices in our pool were retired by
    Microsoft in late April 2026 and now fail with "No audio was received".
    \textbf{Fix:} rebuilt the pool from currently-live multilingual voices,
    verified against \code{edge\_tts.list\_voices()} on 2026-04-25.
  \item[Groq JSON-mode flakiness.] Despite \code{response\_format=\{"type":"json\_object"\}},
    Groq occasionally wrapped the payload in a Markdown fence
    (\code{\textbackslash\textbackslash`\textbackslash`\textbackslash`json...}\code{\textbackslash\textbackslash`\textbackslash`\textbackslash`}) which broke
    \code{json.loads}. \textbf{Fix:} added a defensive extractor that strips
    fences and surrounding prose before parsing, plus a Pydantic
    post-validation layer so a malformed response is degraded to a placeholder
    rather than aborting the whole run.
  \item[WebSocket dropped during long renders.] Cold-start renders that wait
    on Pollinations for 60--90\,s per image would silently drop the progress
    WebSocket on some networks (no events received, but pipeline still
    running). \textbf{Fix:} added a 15-second heartbeat ping from the server
    and a client-side reconnect-with-backoff so the live event log stays
    populated even on flaky links.
  \item[SQLite write contention on parallel edits.] When the web UI submitted
    a \code{/api/edit} while a phase rerun was still writing snapshots, both
    paths competed for \code{state.sqlite} and one would intermittently raise
    \code{database is locked}. \textbf{Fix:} the storage layer now opens
    connections with \code{isolation\_level="IMMEDIATE"} and a 5-second busy
    timeout, serialising writes safely without forcing a global lock.
  \item[Wav2Lip motion-warp dead end.] Before settling on the mouth-region
    blend, we tried "motion-only transfer": run Wav2Lip to extract per-frame
    optical flow, then warp the high-resolution portrait pixels by that flow
    so the original sharp mouth would appear to move with the synthesised
    lip-sync. The warp was technically correct but visually broken --- when
    a mouth opens, the frame contains \emph{new pixels} (teeth, tongue, dark
    interior) that don't exist anywhere in the closed-mouth portrait. Optical
    flow can warp existing pixels but cannot synthesise pixels that aren't
    there; the result looked like the mouth was being smeared shut. Real
    face-reenactment models (FOMM, NeRFace) handle this with a learned prior
    over typical mouth interiors, which was out of scope. \textbf{Resolution:}
    accepted that \emph{some} pixels in the mouth region must come from
    Wav2Lip; minimised the visible footprint with the elliptical mask + CAS
    sharpening described in the Phase 3 section.
  \item[Determinism vs. user expectation on reruns.] The cache-everything,
    deterministic-LLM design (\code{temperature=0.0}, content-hashed tool
    output filenames) gave us byte-identical reproducibility but silently
    converted "Rerun audio / video / story" buttons into no-ops. We
    initially fixed only the route handler, but the bug persisted in the
    edit-agent's \code{audio.rerun} and \code{video.rerun} steps because
    those bypassed the route entirely. \textbf{Fix:} extracted the
    cache-invalidation policy into \code{shared/utils/phase\_cache.py} and
    made every "rerun" code path call it, so the route and the edit agent
    cannot diverge on what "invalidate phase X" means.
\end{description}

% ===================== 5. Testing =====================
\section{Testing}

\textbf{47 passing unit + integration tests} grouped by concern:

\begin{tabularx}{\textwidth}{l X r}
\toprule
\textbf{Module} & \textbf{Coverage} & \textbf{Cases} \\
\midrule
\code{test\_edit\_intent}      & 17 free-text query types $\rightarrow$ correct intent / target / scope, plus offline keyword fallback                                  & 19 \\
\code{test\_state\_manager}    & Append-only log invariant, branching after undo (\code{parent\_version} chain), asset freeze + restore round-trip                     & 4  \\
\code{test\_planner}           & Scene-id parsing from natural language, parameter extraction                                                                          & 4  \\
\code{test\_schemas}           & \code{PipelineState} round-trip, default values                                                                                       & 5  \\
\code{test\_subtitle\_timing}  & Per-line offset accumulation                                                                                                          & 3  \\
\code{test\_mcp\_tools}        & Tool registration, executor retry, auto job\_id injection                                                                             & 5  \\
\code{test\_story\_planner}    & Normalisation of LLM JSON output                                                                                                      & 3  \\
\code{test\_api\_shape}        & Backend route shape via FastAPI \code{TestClient}                                                                                     & 4  \\
\bottomrule
\end{tabularx}

% ===================== 6. Results =====================
\section{Results --- a demo run}

\textbf{Prompt:} "A detective meets a nervous informant in a rainy Tokyo alley
at midnight."  \textbf{Scenes:} 3.  \textbf{Style:} cinematic.

\smallskip
\textbf{Cold-start render time on M4 Air:} approximately 5 minutes, dominated
by Pollinations API wait (4 image fetches at $\sim$30--60 s each). All other
phases combined add \textless\,10 seconds.

\smallskip
\textbf{Output:} a 33-second 1024$\times$576 H.264 + AAC MP4 with: an
establishing wide shot of the alley at the start of each scene, character
close-ups during dialogue (Detective Nakamura's portrait alternating with
Informant Tanaka's), mood-matched BGM under the dialogue, and burned
subtitles. The dialogue actually advances a story (informant divulges
information about the Matsumoto case, detective probes for proof, informant
panics) rather than reading like a weather report.

\smallskip
\textbf{Edits applied:}
\begin{itemize}
  \item \emph{"make scene 1 darker"} \rarr{} BG regenerates with darker prompt
        suffix, \code{eq=brightness=-0.15:saturation=0.9} applied to the scene
        MP4. $\sim$60\,s (Pollinations BG re-fetch).
  \item \emph{"change voice tone to whispered"} \rarr{} per-line TTS
        re-rendered with whisper DSP chain, $\sim$25\,s.
  \item \emph{"speed up scene 2"} \rarr{} \code{atempo=1.5} on scene 2 MP4,
        re-stitch final, $\sim$3\,s.
\end{itemize}

\textbf{Reverts:} two pointer moves --- "Undo last" reverted speed-up; "Restore v3"
returned to the post-darker / pre-whisper state. No new rows were
\emph{created} for the undo (only for the subsequent edits if any).

% ===================== 7. Rubric coverage =====================
\section{Rubric coverage (self-assessment)}

\begin{tabularx}{\textwidth}{l c X}
\toprule
\textbf{Criterion} & \textbf{Weight} & \textbf{Evidence} \\
\midrule
Phase 1 --- Story \& Script    & 15\% & Strict-JSON LLM output, character roster with consistent voice + appearance, validated with \code{test\_story\_planner} \\
Phase 2 --- Audio Generation   & 15\% & Deterministic per-character voice, gender-aware pool, emotion playbook (text cues + DSP), timing manifest accurate to ms \\
Phase 3 --- Video Composition  & 20\% & Per-scene 16:9 native (no upscale), Ken-Burns + close-up flow, A/V sync via timing manifest, xfade transitions, browser-safe \code{yuv420p} + \code{+faststart} \\
Phase 4 --- Web Interface      & 10\% & Live WebSocket progress, scene tile previews, phase re-run, version history with active highlight + lineage tooltip, edit chat \\
Phase 5 --- Edit Agent + Undo  & 20\% & 17 tested intent types, append-only SQLite log + active-pointer + branching via \code{parent\_version}, generic \code{modify\_scene\_visual} for free-text directives \\
Integration / Pipeline         & 10\% & Single Pydantic state object, all phases compose without translation, structured events flow to UI \\
Report \& Presentation         & 10\% & This document + recorded demo \\
\bottomrule
\end{tabularx}

\bigskip
% ===================== 8. Division of work =====================
\section{Division of work}

This was a \textbf{3-person group}. Phase responsibility was distributed so
each member owned roughly equivalent rubric weight, with shared ownership of
the cross-cutting layer (JSON schema, integration testing).

\renewcommand{\arraystretch}{1.4}
\begin{tabularx}{\textwidth}{l l X}
\toprule
\textbf{Member} & \textbf{ID} & \textbf{Primary phase(s) and key responsibilities} \\
\midrule
Hussain Waseem &
22i-0893 &
\textbf{Phase 1 --- Story \& Script.} Pydantic schema design for the
inter-phase JSON contract, Groq LLM integration with strict JSON mode,
story+character generation prompts, character roster validation + portrait
pipeline, edit-intent classifier prompt engineering. \\
Ibtehaj Haider &
22i-0767 &
\textbf{Phase 2 (Audio) + Phase 3 (Video Composition).} Voice pool curation
(edge-tts + ElevenLabs), gender-aware voice picker, emotion prosody map +
DSP playbook, mood-driven BGM (library + synth fallback), dialogue/BGM mixing
with ducking, timing manifest design, Pollinations / fal.\,ai image
generation, FFmpeg primitives (Ken Burns, portrait overlay, color grade),
subtitle burner (Pillow PNG overlays, libass-free), scene + final
compositors with xfade, video agent close-up flow + establishing shots. \\
Muhammad Taha &
22i-0983 &
\textbf{Phase 4 (Web App) + Phase 5 (Edit Agent \& Versioning).} FastAPI +
WebSocket backend, React SPA frontend (prompt $\rightarrow$ progress $\rightarrow$
preview, version history, edit chat), MCP tool registry + executor framework,
LangGraph edit-agent graph (intent classifier + planner + executor),
\code{StateManager} with append-only SQLite log + active-version pointer +
parent-version lineage, undo/restore semantics, end-to-end orchestrator
wiring. \\
\bottomrule
\end{tabularx}
\renewcommand{\arraystretch}{1.25}

\smallskip
Jointly owned by all three members: the shared JSON contract
(\code{shared/schemas/}), integration testing across phase boundaries, the
demo recording, and this report.

% ===================== 9. Future work =====================
\section{Future work}

\begin{itemize}
  \item \textbf{True video generation.} The current stack uses still-image
    Ken-Burns motion, which is stable on CPU and avoids frame-by-frame
    hallucination. A hosted RunwayML / Luma Dream Machine integration could
    plug in as a drop-in replacement for \code{generate\_scene\_background +
    ken\_burns}, behind a feature flag.
  \item \textbf{Better BGM.} The synthesised chord pad is mood-correlated but
    uninspired; expanding the curated CC-licensed library (we have one
    starter set from Serge Quadrado) and adding a local MusicGen path would
    raise production quality meaningfully.
  \item \textbf{Subtitle styling.} Once \code{libass}-enabled \code{ffmpeg} is
    available, ASS subtitles allow per-character typography (italics for
    whispered, larger for shouted) without leaving the FFmpeg pipeline.
  \item \textbf{HD images on a paid tier.} Wiring fal.\,ai schnell behind the
    \code{FAL\_KEY} env var produces native 1920$\times$1080 from FLUX for
    $\sim$\$0.006/image --- worth the few dollars of credit for a final
    submission render.
\end{itemize}

\end{document}
